% 案例式教学方法研究综述
% 编译命令: xelatex main.tex

\documentclass[12pt,a4paper]{ctexart}

% ==================== 宏包导入 ====================
\usepackage{geometry}
\usepackage{graphicx}
\usepackage{booktabs}
\usepackage{longtable}
\usepackage{array}
\usepackage{multirow}
\usepackage{hyperref}
\usepackage{xcolor}
\usepackage{enumitem}
\usepackage{tcolorbox}
\usepackage{fancyhdr}
\usepackage{titlesec}
\usepackage{caption}

% ==================== 页面设置 ====================
\geometry{left=2.5cm, right=2.5cm, top=3cm, bottom=3cm}

% ==================== 超链接设置 ====================
\hypersetup{
    colorlinks=true,
    linkcolor=blue!70!black,
    citecolor=green!50!black,
    urlcolor=blue!80!black
}

% ==================== tcolorbox环境 ====================
\tcbuselibrary{skins, breakable}

\newtcolorbox{keypoint}[1][]{
    colback=blue!5!white,
    colframe=blue!75!black,
    fonttitle=\bfseries,
    title=#1,
    breakable
}

\newtcolorbox{casebox}[1][]{
    colback=green!5!white,
    colframe=green!60!black,
    fonttitle=\bfseries,
    title=#1,
    breakable
}

\newtcolorbox{methodbox}[1][]{
    colback=orange!5!white,
    colframe=orange!70!black,
    fonttitle=\bfseries,
    title=#1,
    breakable
}

% ==================== 文档信息 ====================
\title{\textbf{案例式教学方法研究综述}\\[0.5em]
\large 面向任务规划课程的教学设计参考}
\author{教学研究资料}
\date{2026年2月}

% ==================== 正文开始 ====================
\begin{document}

\maketitle

\begin{abstract}
案例式教学（Case-Based Teaching）是一种以真实或模拟案例为载体，以学生为中心，通过分析、讨论和反思来培养学生解决实际问题能力的教学方法。本文系统综述了案例式教学的起源与发展、国内典型实践、案例资源建设、教学设计方法，并针对任务规划课程提出了案例式教学的设计建议。
\end{abstract}

\tableofcontents
\newpage

% ================================================================
\section{案例式教学概述}
% ================================================================

\subsection{案例式教学的起源}

案例式教学法（Case Method）起源于19世纪末的哈佛大学法学院，由院长Christopher Columbus Langdell首创。1920年代，哈佛商学院将其发展为管理教育的核心方法，形成了著名的"哈佛案例教学法"。

\begin{keypoint}[哈佛案例教学法的核心特征]
\begin{itemize}[leftmargin=2em]
    \item \textbf{真实性}：案例来源于真实的商业情境
    \item \textbf{复杂性}：案例包含多维度的决策因素
    \item \textbf{开放性}：没有标准答案，鼓励多元思考
    \item \textbf{参与性}：学生是课堂讨论的主体
\end{itemize}
\end{keypoint}

哈佛商学院至今已开发超过10,000个教学案例，每年产出约350个新案例，被全球超过80个国家的商学院采用。

\subsection{案例式教学的核心理念}

案例式教学建立在建构主义学习理论基础上，其核心理念包括：

\subsubsection{以学生为中心}

传统教学以教师讲授为主，学生被动接受知识。案例式教学将学生置于决策者的位置，主动参与问题分析和解决方案的制定。教师的角色从"知识传授者"转变为"学习引导者"。

\subsubsection{问题导向学习}

案例式教学以真实问题为驱动，学生需要：
\begin{enumerate}[leftmargin=2em]
    \item 识别案例中的关键问题
    \item 分析问题产生的原因
    \item 收集和评估相关信息
    \item 提出并评价可能的解决方案
    \item 制定实施计划
\end{enumerate}

\subsubsection{经验反思与知识建构}

根据Kolb的经验学习循环理论，有效学习包括四个阶段：具体经验$\rightarrow$反思观察$\rightarrow$抽象概念化$\rightarrow$主动实验。案例式教学通过模拟真实情境，使学生在"安全"的环境中积累经验，并通过讨论和反思将经验转化为可迁移的知识。

\subsection{案例式教学与传统教学的对比}

\begin{table}[htbp]
\centering
\caption{案例式教学与传统教学的对比}
\begin{tabular}{p{3cm}p{5cm}p{5cm}}
\toprule
\textbf{维度} & \textbf{传统教学} & \textbf{案例式教学} \\
\midrule
教学中心 & 以教师为中心 & 以学生为中心 \\
知识来源 & 教材、教师讲授 & 案例分析、讨论交流 \\
学生角色 & 被动接受者 & 主动参与者、决策者 \\
教师角色 & 知识传授者 & 学习引导者、讨论促进者 \\
学习内容 & 抽象理论、概念 & 具体情境、实践问题 \\
评价方式 & 标准化考试 & 参与度、分析能力、方案质量 \\
能力培养 & 记忆、理解 & 分析、综合、评价、创造 \\
\bottomrule
\end{tabular}
\end{table}

\subsection{案例式教学在中国的发展历程}

中国引进案例式教学始于1980年代改革开放时期。发展历程可分为三个阶段：

\begin{enumerate}[leftmargin=2em]
    \item \textbf{引进探索期（1980-2000）}：以翻译、引进国外案例为主，清华大学经管学院、中欧国际工商学院等率先开展案例教学实践。

    \item \textbf{本土化发展期（2000-2015）}：开始大规模开发本土案例，中国管理案例共享中心（CMCC）成立，"全国百篇优秀案例"评选启动。

    \item \textbf{深化创新期（2015至今）}：案例教学向多学科扩展（法学、医学、工程等），数字化案例、虚拟仿真案例兴起，教育部推进新形态教材建设。
\end{enumerate}

% ================================================================
\section{国内典型案例式教学实践}
% ================================================================

\subsection{清华大学经管学院MBA案例教学}

清华大学经济管理学院是中国最早系统开展案例教学的商学院之一。

\subsubsection{中国工商管理案例中心}

清华经管学院于1999年成立中国工商管理案例中心，主要职能包括：
\begin{itemize}[leftmargin=2em]
    \item 开发高质量的本土商业案例
    \item 组织案例教学师资培训
    \item 促进国内外案例资源交流
    \item 推动案例研究方法论发展
\end{itemize}

截至2024年，该中心已开发案例超过600个，涵盖战略管理、组织行为、市场营销、财务管理等领域。

\subsubsection{本土化案例开发经验}

清华经管学院在案例本土化方面积累了丰富经验：

\begin{casebox}[清华经管本土化案例开发原则]
\begin{enumerate}[leftmargin=2em]
    \item \textbf{时代性}：反映中国经济转型、数字化转型等时代特征
    \item \textbf{典型性}：选择具有代表性的中国企业和管理实践
    \item \textbf{理论性}：案例能够承载和验证管理理论
    \item \textbf{教学性}：设计清晰的教学目标和讨论要点
\end{enumerate}
\end{casebox}

\subsection{北京大学法学院诊所式法律教育}

北京大学法学院于2000年引入诊所式法律教育（Clinical Legal Education），这是案例教学在法学领域的重要变体。

\begin{keypoint}[诊所式法律教育的特点]
\begin{itemize}[leftmargin=2em]
    \item 学生在教师指导下代理真实法律案件
    \item 直接接触当事人，了解真实法律需求
    \item 在实践中学习法律技能和职业伦理
    \item 同时提供法律援助服务，具有社会价值
\end{itemize}
\end{keypoint}

北大法学院开设的诊所课程包括：民事法律诊所、刑事法律诊所、劳动法律诊所、妇女权益保护诊所等。这种教学模式培养了学生的实践能力和职业素养。

\subsection{复旦大学医学院PBL教学模式}

复旦大学上海医学院是国内较早引入PBL（Problem-Based Learning，问题导向学习）教学模式的医学院校。

PBL教学的典型流程：
\begin{enumerate}[leftmargin=2em]
    \item \textbf{呈现病例}：教师提供真实或模拟的临床病例
    \item \textbf{界定问题}：学生小组讨论，识别需要解决的医学问题
    \item \textbf{自主学习}：学生根据问题查阅资料，学习相关知识
    \item \textbf{小组讨论}：分享学习成果，整合知识
    \item \textbf{总结评价}：教师引导总结，评估学习效果
\end{enumerate}

\subsection{中国政法大学案例教学体系}

中国政法大学建立了系统的法学案例教学体系，包括：
\begin{itemize}[leftmargin=2em]
    \item \textbf{案例教材}：编写系列法学案例教材
    \item \textbf{案例库}：建设法学案例数据库
    \item \textbf{模拟法庭}：定期组织模拟法庭活动
    \item \textbf{实习基地}：与法院、检察院、律所合作建立实习基地
\end{itemize}

\subsection{中欧国际工商学院案例研究中心}

中欧国际工商学院（CEIBS）案例研究中心是中国领先的商业案例开发机构之一。

\begin{table}[htbp]
\centering
\caption{中欧案例中心主要成果}
\begin{tabular}{ll}
\toprule
\textbf{指标} & \textbf{数据} \\
\midrule
累计开发案例数 & 1000+ \\
年均新增案例 & 50-80个 \\
国际发行案例 & 300+（通过Harvard Publishing等平台）\\
获奖案例 & 多次获得毅伟（Ivey）案例奖 \\
\bottomrule
\end{tabular}
\end{table}

中欧案例的特色在于：关注中国企业的国际化进程、数字化转型、创新创业等前沿议题。

\subsection{工程教育领域案例教学}

\subsubsection{西安交通大学}

西安交通大学在工程教育中推行案例教学，特别是在机械工程、能源动力等专业。主要做法包括：
\begin{itemize}[leftmargin=2em]
    \item 开发工程案例库，涵盖设计、制造、运维等环节
    \item 结合企业实际项目开展案例教学
    \item 将案例教学与课程设计、毕业设计结合
\end{itemize}

\subsubsection{浙江大学}

浙江大学在计算机、软件工程等专业开展案例教学实践：
\begin{itemize}[leftmargin=2em]
    \item 软件工程课程采用真实项目案例
    \item 人工智能课程使用Kaggle竞赛案例
    \item 建设虚拟仿真实验平台，支持案例教学
\end{itemize}

% ================================================================
\section{案例教学资源建设}
% ================================================================

\subsection{中国管理案例共享中心（CMCC）}

中国管理案例共享中心（China Management Case-sharing Center, CMCC）由全国MBA教育指导委员会牵头，于2007年正式成立，是中国最大的管理案例资源平台。

\subsubsection{案例库建设}

CMCC案例库的建设采用"征集-评审-入库-共享"的模式：

\begin{enumerate}[leftmargin=2em]
    \item \textbf{案例征集}：面向全国商学院征集优秀案例
    \item \textbf{双盲评审}：邀请专家进行匿名评审
    \item \textbf{入库管理}：通过评审的案例统一入库
    \item \textbf{共享使用}：成员单位可免费使用库中案例
\end{enumerate}

截至2024年，CMCC案例库收录案例超过5000个。

\subsubsection{案例入库标准}

\begin{methodbox}[CMCC案例入库评审标准]
\begin{enumerate}[leftmargin=2em]
    \item \textbf{真实性}（20\%）：基于真实企业和事件
    \item \textbf{教学性}（25\%）：明确的教学目标，丰富的讨论空间
    \item \textbf{理论性}（20\%）：能够承载相关管理理论
    \item \textbf{写作质量}（20\%）：结构清晰，叙述流畅
    \item \textbf{配套材料}（15\%）：教学笔记、课件、数据附件等
\end{enumerate}
\end{methodbox}

\subsection{教育部新形态教材建设要求}

教育部近年来大力推进新形态教材建设，对案例式教材提出了明确要求。

\subsubsection{数字化资源融合}

新形态教材强调纸质教材与数字资源的深度融合：
\begin{itemize}[leftmargin=2em]
    \item 二维码链接扩展资源（视频、数据、仿真）
    \item 在线案例讨论平台
    \item 虚拟仿真实验
    \item 智能题库与自适应学习
\end{itemize}

\subsubsection{案例式教材编写规范}

\begin{keypoint}[案例式教材编写要点]
\begin{enumerate}[leftmargin=2em]
    \item 每章以案例导入，激发学习兴趣
    \item 理论讲解与案例分析交叉进行
    \item 设置渐进式案例，由浅入深
    \item 配套讨论题和延伸阅读
    \item 提供教师版教学笔记
\end{enumerate}
\end{keypoint}

\subsection{优秀案例的评选标准}

\subsubsection{全国百篇优秀案例}

全国MBA教育指导委员会每年评选"全国百篇优秀管理案例"，这是中国管理案例领域的最高荣誉。评选标准包括：

\begin{table}[htbp]
\centering
\caption{全国百篇优秀案例评选维度}
\begin{tabular}{p{3cm}p{8cm}}
\toprule
\textbf{评选维度} & \textbf{具体要求} \\
\midrule
案例质量 & 真实性、典型性、时效性、深度 \\
教学适用性 & 教学目标明确、讨论空间充分、难度适中 \\
理论贡献 & 能够验证或发展管理理论 \\
写作规范 & 结构完整、语言流畅、数据翔实 \\
实际使用 & 已在教学中使用并取得良好效果 \\
\bottomrule
\end{tabular}
\end{table}

\subsubsection{毅伟案例奖}

毅伟案例奖（Ivey/CEIBS Case Award）由加拿大毅伟商学院与中欧国际工商学院联合设立，旨在表彰亚洲地区的优秀商业案例。获奖案例通常具有以下特点：
\begin{itemize}[leftmargin=2em]
    \item 聚焦亚洲（特别是中国）商业实践
    \item 具有国际视野和普适性教学价值
    \item 写作质量达到国际出版标准
\end{itemize}

\subsection{案例开发流程}

\subsubsection{选题与调研}

案例选题是开发的第一步，需要考虑：
\begin{enumerate}[leftmargin=2em]
    \item \textbf{教学需求}：课程需要什么样的案例？
    \item \textbf{理论承载}：案例能够说明什么理论？
    \item \textbf{可获取性}：能否获得足够的信息和数据？
    \item \textbf{时效性}：案例是否反映当前商业环境？
\end{enumerate}

调研方式包括：企业访谈、公开资料收集、行业报告分析等。

\subsubsection{案例撰写}

\begin{methodbox}[案例文本结构]
\begin{enumerate}[leftmargin=2em]
    \item \textbf{开篇}：设置悬念，引出核心问题
    \item \textbf{背景}：企业/行业/人物背景介绍
    \item \textbf{正文}：按时间或逻辑顺序叙述事件经过
    \item \textbf{决策点}：明确决策者面临的选择
    \item \textbf{附录}：财务数据、组织架构、补充材料
\end{enumerate}
\end{methodbox}

\subsubsection{教学笔记编写}

教学笔记（Teaching Note）是供教师使用的配套文档，通常包括：
\begin{itemize}[leftmargin=2em]
    \item 案例概要与教学目标
    \item 建议的课程安排（时间分配）
    \item 讨论题及参考答案
    \item 板书设计建议
    \item 可能的学生反应及应对策略
    \item 后续发展（Epilogue）
\end{itemize}

\subsubsection{案例测试与修订}

案例在正式发布前应进行测试：
\begin{enumerate}[leftmargin=2em]
    \item 小范围试用：在一两个班级试用
    \item 收集反馈：学生反馈、教师反馈
    \item 修订完善：根据反馈修改案例和教学笔记
    \item 同行评审：邀请同行专家评审
\end{enumerate}

% ================================================================
\section{案例式教学设计方法}
% ================================================================

\subsection{三阶段教学模式}

基于国内成熟实践，案例式教学通常采用"三阶段"模式：

\subsubsection{个人学习阶段}

\begin{itemize}[leftmargin=2em]
    \item \textbf{课前阅读}：学生独立阅读案例材料（通常需要1-2小时）
    \item \textbf{问题思考}：完成教师布置的预习思考题
    \item \textbf{资料查阅}：根据需要查阅相关背景资料
    \item \textbf{形成观点}：对案例问题形成初步判断
\end{itemize}

\subsubsection{小组讨论阶段}

\begin{itemize}[leftmargin=2em]
    \item \textbf{组内分享}：每位成员分享自己的分析和观点
    \item \textbf{观点碰撞}：讨论不同观点的合理性
    \item \textbf{达成共识}：形成小组的主要结论
    \item \textbf{准备汇报}：准备课堂汇报材料
\end{itemize}

小组通常由4-6人组成，讨论时间30-60分钟。

\subsubsection{课堂讨论阶段}

\begin{methodbox}[课堂讨论教师引导技巧]
\begin{enumerate}[leftmargin=2em]
    \item \textbf{开场提问}：用开放性问题启动讨论
    \item \textbf{追问深化}："为什么这样认为？""还有其他可能吗？"
    \item \textbf{对比分析}：引导不同观点的对比
    \item \textbf{联系理论}：适时引入相关理论框架
    \item \textbf{总结升华}：提炼讨论要点，升华到一般性结论
\end{enumerate}
\end{methodbox}

\subsection{案例难度阶梯设计}

为适应不同学习阶段，同一主题的案例可设计多个难度版本：

\begin{table}[htbp]
\centering
\caption{案例难度阶梯设计}
\begin{tabular}{p{2cm}p{4cm}p{6cm}}
\toprule
\textbf{版本} & \textbf{特点} & \textbf{适用场景} \\
\midrule
简化版 & 问题明确，信息充分，干扰少 & 初学者入门，概念理解 \\
标准版 & 信息适度复杂，有一定干扰 & 正常教学，能力训练 \\
进阶版 & 信息不完整，多目标冲突 & 高年级学生，综合应用 \\
\bottomrule
\end{tabular}
\end{table}

\subsection{讨论题设计原则}

好的讨论题应满足以下原则：

\begin{enumerate}[leftmargin=2em]
    \item \textbf{开放性}：没有唯一正确答案
    \item \textbf{挑战性}：需要深入思考才能回答
    \item \textbf{关联性}：与教学目标紧密相关
    \item \textbf{渐进性}：由浅入深，逐步深化
    \item \textbf{可讨论性}：能够引发不同观点的碰撞
\end{enumerate}

\begin{casebox}[讨论题设计示例]
\textbf{描述性问题}：案例中的主要问题是什么？\\
\textbf{分析性问题}：造成这个问题的根本原因是什么？\\
\textbf{评价性问题}：你如何评价决策者的做法？\\
\textbf{创造性问题}：如果你是决策者，你会怎么做？
\end{casebox}

\subsection{案例教学效果评估}

案例教学的评估应关注过程和能力，而非仅仅是知识记忆：

\begin{table}[htbp]
\centering
\caption{案例教学评估维度}
\begin{tabular}{p{3cm}p{3cm}p{5cm}}
\toprule
\textbf{评估维度} & \textbf{权重建议} & \textbf{评估方式} \\
\midrule
课堂参与 & 20-30\% & 发言质量、讨论贡献 \\
书面分析 & 30-40\% & 案例分析报告 \\
小组表现 & 15-25\% & 小组讨论、汇报 \\
期末考核 & 20-30\% & 案例分析考试 \\
\bottomrule
\end{tabular}
\end{table}

\subsection{常见问题与应对策略}

\begin{table}[htbp]
\centering
\caption{案例教学常见问题及应对}
\begin{tabular}{p{4cm}p{7cm}}
\toprule
\textbf{问题} & \textbf{应对策略} \\
\midrule
学生准备不充分 & 设置预习检查，将预习纳入成绩 \\
讨论冷场 & 使用点名提问、小组轮流发言 \\
讨论跑题 & 教师适时引导，回归核心问题 \\
少数人主导 & 鼓励安静学生发言，控制活跃学生 \\
时间控制困难 & 提前规划，设置时间节点 \\
理论联系不足 & 讨论后专门环节进行理论总结 \\
\bottomrule
\end{tabular}
\end{table}

% ================================================================
\section{面向任务规划课程的案例教学设计}
% ================================================================

\subsection{任务规划案例的特点}

任务规划课程的案例具有以下特殊性：

\subsubsection{问题复杂性}

任务规划问题通常涉及：
\begin{itemize}[leftmargin=2em]
    \item 多目标优化（效率、成本、安全等）
    \item 复杂约束（时间、资源、物理）
    \item 不确定性因素（环境、敌情、设备故障）
    \item 多智能体协调（分布式决策、信息共享）
\end{itemize}

\subsubsection{算法多样性}

同一问题可能适用多种算法方法：
\begin{itemize}[leftmargin=2em]
    \item 经典规划方法（STRIPS/PDDL、HTN）
    \item 约束求解方法（CSP、SAT、SMT）
    \item 优化方法（启发式搜索、元启发式）
    \item 学习方法（强化学习、模仿学习）
    \item 博弈方法（博弈规划、对抗规划）
\end{itemize}

这种多样性为案例教学提供了丰富的讨论空间。

\subsubsection{应用场景真实性}

任务规划案例来源于真实应用场景：
\begin{itemize}[leftmargin=2em]
    \item 航天领域：卫星观测调度、星座部署
    \item 军事领域：作战任务规划、无人机协同
    \item 物流领域：车辆路径规划、仓储调度
    \item 机器人领域：运动规划、任务分配
\end{itemize}

\subsection{案例选取原则}

\subsubsection{代表性与典型性}

选择能够代表任务规划核心问题的案例：
\begin{itemize}[leftmargin=2em]
    \item 覆盖主要问题类型（调度、路径、分配）
    \item 涉及关键技术难点（约束处理、不确定性、协调）
    \item 反映实际应用需求
\end{itemize}

\subsubsection{可拓展性}

案例应具有可拓展性，便于设计不同难度版本：
\begin{itemize}[leftmargin=2em]
    \item 从简单到复杂的问题递进
    \item 从确定到不确定的环境变化
    \item 从单体到多体的规模扩展
    \item 从合作到对抗的关系演变
\end{itemize}

\subsubsection{教学适切性}

案例应适合本科生的知识水平和学习能力：
\begin{itemize}[leftmargin=2em]
    \item 问题规模适中，便于分析
    \item 数据完整，支持算法实现
    \item 有标准解或已知上界，便于评估
\end{itemize}

\subsection{贯穿式案例设计}

\subsubsection{案例版本递进}

通过同一案例的不同版本贯穿多个章节，形成知识体系的连贯性：

\begin{casebox}[案例版本递进示例：多星观测任务]
\begin{itemize}[leftmargin=2em]
    \item \textbf{v1（第2章）}：单星观测，确定性环境 $\rightarrow$ PDDL建模
    \item \textbf{v2（第3章）}：增加时间窗口、能量约束 $\rightarrow$ HTN分解、约束求解
    \item \textbf{v3（第4章）}：引入云层遮挡不确定性 $\rightarrow$ MDP/POMDP
    \item \textbf{v4（第5章）}：多星协同覆盖 $\rightarrow$ 分布式规划、MARL
    \item \textbf{v5（第6章）}：对抗性观测（敌方干扰）$\rightarrow$ 博弈规划
\end{itemize}
\end{casebox}

\subsubsection{章节间逻辑衔接}

每个新版本引入时，需要说明：
\begin{enumerate}[leftmargin=2em]
    \item \textbf{为什么需要新方法}：前一版本的方法在新场景下的局限性
    \item \textbf{新方法解决什么问题}：新方法的针对性和优势
    \item \textbf{方法之间的关系}：是替代还是补充？能否组合使用？
\end{enumerate}

\subsection{本课程案例教学实施建议}

\begin{methodbox}[任务规划课程案例教学实施建议]
\begin{enumerate}[leftmargin=2em]
    \item \textbf{案例导入}：每章以案例场景开始，提出待解决的规划问题
    \item \textbf{算法讲解}：介绍解决该问题的算法原理
    \item \textbf{案例分析}：用算法分析案例，展示建模和求解过程
    \item \textbf{编程实践}：学生动手实现算法，在案例数据上运行
    \item \textbf{讨论比较}：讨论不同方法的优劣，分析适用场景
    \item \textbf{案例演化}：引出案例的下一版本，为后续章节做铺垫
\end{enumerate}
\end{methodbox}

% ================================================================
\section{结论与展望}
% ================================================================

\subsection{案例式教学的优势与挑战}

\subsubsection{主要优势}

\begin{enumerate}[leftmargin=2em]
    \item \textbf{提升学习动机}：真实案例激发学习兴趣
    \item \textbf{培养实践能力}：在模拟情境中锻炼决策能力
    \item \textbf{促进深度学习}：讨论和反思促进知识内化
    \item \textbf{发展综合素质}：沟通、协作、批判性思维
\end{enumerate}

\subsubsection{主要挑战}

\begin{enumerate}[leftmargin=2em]
    \item \textbf{案例开发}：高质量案例开发耗时耗力
    \item \textbf{教师能力}：需要教师具备引导讨论的能力
    \item \textbf{学生适应}：部分学生不适应主动学习方式
    \item \textbf{评估困难}：能力评估比知识测试更复杂
\end{enumerate}

\subsection{数字化转型中的案例教学}

数字化技术为案例教学带来新机遇：

\begin{itemize}[leftmargin=2em]
    \item \textbf{虚拟仿真}：构建可交互的案例环境
    \item \textbf{数据驱动}：使用真实数据增强案例真实性
    \item \textbf{在线协作}：支持异步和远程的案例讨论
    \item \textbf{智能辅助}：AI辅助案例分析和评估
\end{itemize}

\subsection{未来发展趋势}

\begin{enumerate}[leftmargin=2em]
    \item \textbf{跨学科案例}：打破学科边界，开发综合性案例
    \item \textbf{动态案例}：案例随时间演化，反映最新发展
    \item \textbf{个性化学习}：根据学生水平推荐适合的案例版本
    \item \textbf{国际化合作}：中外案例资源共享与联合开发
\end{enumerate}

% ================================================================
% 参考文献
% ================================================================
\newpage
\section*{参考文献}
\addcontentsline{toc}{section}{参考文献}

\begin{enumerate}[leftmargin=2em, label={[\arabic*]}]
    \item 苏敬勤, 崔淼. 案例研究方法：理论与实践[M]. 北京: 清华大学出版社, 2019.
    \item 毛基业, 李晓燕. 案例研究的"术"与"道"的反思[J]. 管理世界, 2021(3): 248-261.
    \item 全国MBA教育指导委员会. 中国管理案例共享中心案例入库标准[S]. 2020.
    \item 教育部. 新形态教材建设指南[Z]. 2021.
    \item Barnes, L.B., Christensen, C.R., Hansen, A.J. Teaching and the Case Method[M]. Harvard Business School Press, 1994.
    \item Kolb, D.A. Experiential Learning: Experience as the Source of Learning and Development[M]. Prentice Hall, 1984.
    \item Erskine, J.A., Leenders, M.R., Mauffette-Leenders, L.A. Teaching with Cases[M]. Ivey Publishing, 2003.
    \item 清华大学经济管理学院中国工商管理案例中心. 案例开发与教学指南[M]. 北京: 清华大学出版社, 2018.
    \item 王玉, 徐淑英. 管理学研究中的案例研究方法[J]. 管理世界, 2007(4): 146-155.
    \item 陈春花, 刘祯. 中国情境下的管理研究[J]. 管理学报, 2017(1): 1-10.
    \item 杨斌. MBA案例教学的本土化实践[J]. 学位与研究生教育, 2015(6): 1-6.
    \item 梁正, 陈劲. 工程教育中的案例教学[J]. 高等工程教育研究, 2019(2): 78-84.
\end{enumerate}

% ================================================================
% 附录
% ================================================================
\newpage
\appendix

\section{国内主要案例资源平台}

\begin{table}[htbp]
\centering
\caption{国内主要案例资源平台}
\begin{tabular}{p{4cm}p{4cm}p{4cm}}
\toprule
\textbf{平台名称} & \textbf{主办单位} & \textbf{网址} \\
\midrule
中国管理案例共享中心 & 全国MBA教指委 & cmcc.shufe.edu.cn \\
清华案例中心 & 清华大学经管学院 & case.sem.tsinghua.edu.cn \\
中欧案例中心 & 中欧国际工商学院 & casecenter.ceibs.edu \\
北大案例研究中心 & 北京大学光华管理学院 & case.gsm.pku.edu.cn \\
毅伟案例库 & Ivey商学院 & iveycases.com \\
哈佛案例库 & Harvard Business Publishing & hbsp.harvard.edu \\
\bottomrule
\end{tabular}
\end{table}

\section{案例模板}

\subsection{案例正文模板}

\begin{verbatim}
案例标题

【导言】
设置场景，引出核心问题（100-200字）

【背景】
1. 行业背景
2. 企业/组织背景
3. 关键人物介绍

【正文】
按时间或逻辑顺序叙述事件发展
- 事件一
- 事件二
- ...
- 决策点

【附录】
- 财务数据/技术参数
- 组织架构图
- 时间线
- 其他补充材料
\end{verbatim}

\subsection{教学笔记模板}

\begin{verbatim}
教学笔记

【案例概要】（200字以内）

【教学目标】
1. 知识目标
2. 能力目标
3. 态度目标

【适用课程与学时】

【讨论题】
1. 基础问题
2. 分析问题
3. 综合问题

【建议板书】

【讨论要点与参考答案】

【后续发展】（如有）

【延伸阅读】
\end{verbatim}

\end{document}
